\documentclass[../../main.tex]{subfiles}
\begin{document}
\section{映射到子流形上的限制}\label{section:ch05-限制映射到子流形}

设 $S\subseteq M$ 是浸入子流形。由于 $S$ 上的拓扑可能与子空间拓扑不同,把光滑映射限制到 $S$ 上并非总是一件简单的事。在本节,我们考察把映射的定义域或陪域限制到光滑子流形上的含义。

我们先有一个引理。回忆第一章, $M$ 的子集称为闭测地的,如果它是某个光滑函数 $f:M\to\mathbb{R}$ 的(闭)零点集,该函数满足对 $p\in Z(f)$ 有 $\mathrm{d}f_p\neq 0$。由闭集包含引理(\cref{lemma:ch01-closed-set-collar}),每个闭测地子集是某光滑函数的零点集,其零点集恰为该子集。

\begin{lemma}[Extension Lemma]\label{lemma:ch05-extension}
设 $S\subseteq M$ 是嵌入子流形, $U$ 是 $S$ 在 $M$ 中的任一邻域。对任一光滑函数 $f:S\to\mathbb{R}$,存在 $M$ 上的光滑函数 $\widetilde{f}$,使 $\widetilde{f}|_S=f$;若 $S$ 闭且 $f$ 的支集含于 $S$ 的某闭子集 $A$ 中,则可选取 $\widetilde{f}$ 使其支集含于 $U$。
\end{lemma}

\begin{proof}
我们先证明 $S$ 的每一点 $p$ 都有邻域 $V_p\subseteq M$,使 $f$ 可延拓为 $V_p$ 上的光滑函数。由于 $S$ 是嵌入子流形,我们可选取以 $p$ 为中心的 $S$ 在 $M$ 中的切片坐标卡 $(V_p,\varphi)$,使 $S\cap V_p$ 由 $x^{k+1}=\cdots=x^n=0$ 给出。在该坐标下, $f$ 是 $x^1,\dots,x^k$ 的光滑函数,而它显然可延拓为 $V_p$ 上 $x^1,\dots,x^n$ 的光滑函数,只需令它与 $x^{k+1},\dots,x^n$ 无关即可。于是 $\widetilde{f}$ 在覆盖 $S$ 的开集族 $\{V_p:p\in S\}$ 上有定义。

由闭集包含引理,存在闭测地子集 $Z\subseteq M\setminus S$。若 $S$ 闭,我们可以取 $U$ 使它包含 $f$ 的支集;若 $S$ 不闭,我们就忽略这一要求。由\cref{theorem:ch01-partition-unity},存在光滑单位分解 $\{\psi_p,\psi_0\}$ 从属开覆盖 $\{V_p:p\in S\}\cup\{M\setminus Z\}$。定义 $\widetilde{f}:M\to\mathbb{R}$ 为
\begin{align*}
\widetilde{f}(x)=\sum_p \psi_p(x)\widetilde{f}_p(x),
\end{align*}
其中 $\widetilde{f}_p$ 是 $f$ 在 $V_p$ 上的延拓。这个和局部有限,且每个 $\widetilde{f}_p$ 在 $V_p$ 上光滑,故 $\widetilde{f}$ 在 $M$ 上光滑。若 $x\in S$,则对每个 $p$, $\psi_0(x)\widetilde{f}_p(x)=0$,因为 $\psi_0$ 的支集与 $S$ 不交;而对 $x\in V_p$, $\sum_p\psi_p(x)=1$,故 $\widetilde{f}|_S=f$。

最后,若 $S$ 闭且 $f$ 的支集含于闭子集 $A\subseteq S$,上面构造的 $\widetilde{f}$ 的支集必含于 $U\cup(M\setminus S)=U$。
\end{proof}

注意,上一引理中的延拓 $\widetilde{f}$ 在 $S$ 上取与 $f$ 相同的值,但在 $S$ 附近不一定是 $f$ 的好逼近。

\begin{theorem}[Restricting the Domain]\label{theorem:ch05-restrict-domain}
设 $M$ 是带边或不带边的光滑流形, $S\subseteq M$ 是浸入子流形, $F:M\to N$ 是光滑映射。则限制映射 $F|_S:S\to N$(赋予 $S$ 由其浸入子流形结构诱导的拓扑与光滑结构)是光滑的。
\end{theorem}

\begin{proof}
这是一个局部问题:对每个 $p\in S$,只需证明 $F|_S$ 在 $p$ 的邻域上光滑。由\cref{proposition:ch05-immersed-slice-local},存在 $M$ 的坐标卡 $(U,\varphi)$ 使 $S\cap U$ 是 $U$ 的单个切片,在切片坐标下为 $x^{k+1}=\cdots=x^n=0$。由于限制到开子集不影响光滑性,我们只需把 $F$ 限制到 $U\cap S$ 上。在该切片坐标下, $F$ 的坐标表示是 $\widehat{F}(x^1,\dots,x^n)$,而 $F|_{S\cap U}$ 的坐标表示是 $\widehat{F}(x^1,\dots,x^k,0,\dots,0)$,它是光滑的。
\end{proof}

下面的定理在把光滑映射的陪域限制为子流形时,给出光滑性所受影响的完整描述。

\begin{theorem}[Restricting the Codomain]\label{theorem:ch05-restrict-codomain}
设 $S\subseteq M$ 是浸入子流形。
\begin{enumerate}[(a)]
\item 若 $F:N\to M$ 是光滑映射,且其像实际含于 $S$ 中,则 $F:N\to S$(看作到 $S$ 本身的映射)是光滑的,当且仅当它是到 $M$ 的光滑映射。
\item 若 $F:M\to N$ 是光滑映射,则限制映射 $F|_S:S\to N$ 是光滑的。
\end{enumerate}
\end{theorem}

\begin{proof}
若 $F:N\to S$ 是到 $S$ 的光滑映射,则它与光滑包含映射 $\iota:S\hookrightarrow M$ 的复合是到 $M$ 的光滑映射 $F:N\to M$。反之,设 $F:N\to M$ 光滑且 $F(N)\subseteq S$。令 $p\in N$ 任意,并记 $q=F(p)\in S$。由于 $\iota:S\hookrightarrow M$ 是浸入,局部浸入定理(\cref{theorem:ch04-local-embedding})保证存在 $S$ 的邻域 $U$,使 $\iota|_U:U\to M$ 是嵌入。必要时用原像替换 $N$,可设 $F(N)\subseteq U$。然后选取 $U$ 在 $M$ 中的切片坐标卡 $(V,\varphi)$,使 $U$ 对应于 $x^{k+1}=\cdots=x^n=0$。在这些坐标下, $F:N\to M$ 的坐标表示为 $\widehat{F}(x)=\bigl(\widehat{F}^1(x),\dots,\widehat{F}^n(x)\bigr)$,而由于其像含于 $U$ 中,必有 $\widehat{F}^{k+1}(x)=\cdots=\widehat{F}^n(x)=0$。因此 $F:N\to U$ 的坐标表示只是 $\bigl(\widehat{F}^1(x),\dots,\widehat{F}^k(x)\bigr)$,它是光滑的。这就证明了 (a)。

(b) 是 (a) 的一个特例,因为 $F|_S$ 等于 $F$ 接着 $S$ 到 $M$ 的包含映射。
\end{proof}

\begin{figure}[ht]
\centering
\includegraphics[width=0.8\textwidth]{fig-05-10.png}
\caption{限制陪域}\label{figure:ch05-restrict-codomain}
\end{figure}

\begin{figure}[ht]
\centering
\includegraphics[width=0.8\textwidth]{fig-05-11.png}
\caption{定理 5.29 的证明}\label{figure:ch05-proof-thm-529}
\end{figure}

\begin{example}[A Smooth Map into a Dense Curve]\label{example:ch05-smooth-into-dense}
设 $\gamma:\mathbb{R}\to\mathbb{T}^2$ 是\cref{example:ch04-dense-curve} 中的单射浸入。由于像 $\gamma(\mathbb{R})$ 在 $\mathbb{T}^2$ 中稠密,我们可问:是否存在任何连续映射 $F:S^1\to\mathbb{T}^2$,其像含于 $\gamma(\mathbb{R})$ 中,并作为到 $\mathbb{T}^2$ 的映射光滑?

对这个问题,答案是肯定的,但连续映射 $F$ 到 $\gamma(\mathbb{R})$ 上的提升 $S^1\to\mathbb{R}$ 几乎从不光滑。设 $\alpha(t)=(e^{2\pi\mathrm{i}t},e^{2\pi\mathrm{i}\alpha t})$,其中 $\alpha$ 是无理数。若 $F:S^1\to\mathbb{T}^2$ 连续且像含于 $\gamma(\mathbb{R})$ 中,则由单值性定理(\cref{proposition:appA-monodromy}), $F$ 可提升为连续映射 $\widetilde{F}:S^1\to\mathbb{R}$,满足 $\gamma\circ\widetilde{F}=F$。若 $\widetilde{F}$ 光滑,则 $F$ 也光滑,因为它是光滑映射之复合;但反过来, $F$ 到 $\mathbb{T}^2$ 的光滑性并不保证 $\widetilde{F}$ 到 $\mathbb{R}$ 的光滑性。事实上,例如,若 $\widetilde{F}:S^1\to\mathbb{R}$ 是连续但非光滑的函数,则 $F=\gamma\circ\widetilde{F}$ 未必光滑;但若 $\widetilde{F}(S^1)$ 含于 $\gamma$ 的定义域中的某区间 $I$ 内,使 $\gamma|_I$ 是嵌入,则 $\widetilde{F}=(\gamma|_I)^{-1}\circ F$ 光滑。
\end{example}

在本节结束之前,我们证明关于光滑子流形光滑结构唯一性的一个重要结果。

\begin{theorem}[Uniqueness of Smooth Structures on Submanifolds]\label{theorem:ch05-uniqueness}
设 $M$ 是带边或不带边的光滑流形, $S\subseteq M$ 是子集。
\begin{enumerate}[(a)]
\item 若 $S$ 是 $M$ 的嵌入子流形,则在 $S$ 上使 $S$ 成为 $M$ 的嵌入子流形的光滑结构唯一。
\item 若 $S$ 是 $M$ 的浸入子流形,则在 $S$ 上使 $S$ 成为 $M$ 的浸入子流形的拓扑唯一,但光滑结构不一定唯一。
\end{enumerate}
\end{theorem}

\begin{proof}
先设 $S\subseteq M$ 是嵌入子流形,并设 $S$ 上有两个光滑结构使它成为 $M$ 的嵌入子流形。记相应的光滑流形分别为 $S_1$ 与 $S_2$。由于在两种情形下 $S$ 上的拓扑都是子空间拓扑,映射 $\operatorname{Id}_S:S_1\to S_2$ 是同胚。为证它是微分同胚,我们必须证明它及其逆都是光滑的。我们将证明 $\operatorname{Id}_S:S_1\to S_2$ 光滑;另一方向同理。

设 $\iota_j:S_j\hookrightarrow M$ 为包含映射。由于 $\iota_2:S_2\to M$ 是光滑映射, \cref{theorem:ch05-restrict-codomain}(b) 表明 $\iota_2|_{S_1}:S_1\to M$ 是光滑的。但 $\iota_2|_{S_1}=\iota_2\circ\operatorname{Id}_S$。由于 $\iota_2$ 是浸入, \cref{theorem:ch05-restrict-codomain}(a) 表明 $\operatorname{Id}_S:S_1\to S_2$ 光滑。这就证明了 (a)。

为证 (b),先注意到若 $S$ 是 $M$ 的浸入子流形,则 $S$ 上的拓扑必须使包含映射成为浸入,因而局部嵌入,故它由浸入子流形结构唯一确定。(它就是使包含映射为浸入的最粗拓扑,而这是浸入子流形定义的一部分。)然而,一般没有唯一的光滑结构使 $S$ 成为浸入子流形: \cref{example:ch05-immersed-examples}(c) 中的 $S$ 可被赋予至少两个不同的光滑结构,每个都使包含映射成为单射浸入。
\end{proof}

\end{document}